% ============================================================
% SOLUTION — PART VI: CASE STUDY
%
% AI INSTRUCTION:
% - Provide a model answer for EVERY scenario, not just one ---
%   different students chose differently
% - Open with the grading rubric so the student can self-mark
% - Diagrams: describe the expected diagram in words or draw it with tikz
% ============================================================

\section*{Part VI: Case Study --- Solution}

\begin{keypoint}[Rubric]
\begin{itemize}
    \item Correctness of concepts --- N pts
    \item Justification and reasoning --- N pts
    \item Clarity of presentation --- N pts
\end{itemize}
Any one scenario answered well earns full marks; the scenarios are equivalent
in difficulty.
\end{keypoint}

\subsection*{Scenario A --- Title of the First Scenario}

\textbf{(a)}
\begin{ans}
Walk through the events one at a time. For each: which component detects it,
which decides, which acts.
\end{ans}

\textbf{(b)}
\begin{ans}
Describe the expected end state precisely (what runs where, how many of each).
A labelled sketch is acceptable; state what must appear in it to earn the marks.
\end{ans}

\textbf{(c)}
\begin{ans}
The extended-failure analysis, in terms of the architecture: what keeps
working, what stops, and why.
\end{ans}

\subsection*{Scenario B --- Title of the Second Scenario}

\begin{ans}
Model answer for (a), (b), (c).
\end{ans}

\subsection*{Scenario C --- Title of the Third Scenario}

\begin{ans}
Model answer for (a), (b), (c).
\end{ans}

\begin{caution}
Answering more than one scenario scores zero for the part --- this is stated
on the exam paper and is enforced.
\end{caution}
